% Einheit 2 von 5 – Winkel an Geradenkreuzungen und Parallelen (bank/winkel-dreiecke/e2.jsonl, zusammenbau v0.1)
%% TODO Zweigzeile Teil 1: was hier gelernt wird (Satz in Schülersprache aus dem Einheitstitel) – Einheit 2
\einheitenkopf[e2]{Einheit 2 von 5 · Winkel an Geradenkreuzungen und Parallelen}
\zweigzeile{ab Kl. 6, je nach Buch bis Kl. 7 (am Gymnasium ab Kl. 5) · P10 oft}
\setcounter{aufgabe}{11}

%% TODO Titel: Ich-kann-Satz fehlt in der Bank (Platzhalter: Kettenname) – Nr. 12
%% TODO Anweisung: ein Satz über der ersten Teilaufgabe fehlt in der Bank – Nr. 12
\begin{aufgabe}{Sind die Geraden parallel?}
\begin{teile}
\teil Angabe: Die Geraden g und h sind parallel; die Gerade k schneidet beide. Darfst du g und h als parallel ansehen? Markiere Parallelen mit Pfeilen oder kreuze an: \\ \kreuz{als parallel angegeben} \\ \kreuz{nicht angegeben}

\parallelenpaar{58}{}{}{}{}
\end{teile}
\end{aufgabe}

%% TODO Titel: Ich-kann-Satz fehlt in der Bank (Platzhalter: Kettenname) – Nr. 13
%% TODO Anweisung: ein Satz über der ersten Teilaufgabe fehlt in der Bank – Nr. 13
\begin{aufgabe}{Winkel an Geraden}
%% TODO \rechenplatz{n}: Schrittzahl des Grundfalls fehlt in der Bank (nur bei fester Schreibform, 2.3 b)
\begin{teile}
\teil Wie liegen α und γ an der Kreuzung zueinander? \\ \kreuz{Scheitelwinkel} \\ \kreuz{Nebenwinkel} \\ \kreuz{Stufenwinkel} \\ \kreuz{Wechselwinkel}

\geradenkreuzung{55}{\alpha}{\beta}{\gamma}{\delta}
\teil An einer Geradenkreuzung ist $\alpha = 38^\circ$. Wie groß ist der Scheitelwinkel von α?

\geradenkreuzung{38}{\alpha}{}{}{}
\leerfeld °
\teil An einer Geradenkreuzung ist $\alpha = 44^\circ$. Wie groß ist ein Nebenwinkel von α?

\geradenkreuzung{44}{\alpha}{}{}{}
\leerfeld °
\teil An einer Geradenkreuzung ist $\alpha = 72^\circ$; β und δ liegen neben α, γ liegt α gegenüber. Gib β, γ und δ an.

\geradenkreuzung{72}{\alpha}{\beta}{\gamma}{\delta}
β = \leerfeld °, γ = \leerfeld °, δ = \leerfeld °
\teil Die Geraden g und h schneiden sich; zwischen ihnen liegt $\alpha = 64^\circ$. Eine dritte Gerade durch den Schnittpunkt teilt den Scheitelwinkel von α in zwei Teile, der untere misst $23^\circ$. Wie groß ist der obere Teil β? \hfill (P10 2014 OS) \\ β = \leerfeld °
\teil $g \parallel h$, g liegt oben, h unten, die Gerade k schneidet beide. An g liegt oberhalb von g rechts von k ein Winkel von $58^\circ$. Wie groß ist der Winkel an h oberhalb von h rechts von k?

\parallelenpaar{58}{}{}{}{}
\leerfeld °
\teil $g \parallel h$, g liegt oben, h unten, die Gerade k schneidet beide. An g liegt unterhalb von g links von k ein Winkel von $49^\circ$. Wie groß ist der Winkel an h oberhalb von h rechts von k? \hfill (P10 2020 OS) \\

\parallelenpaar{49}{}{}{}{}
\leerfeld °
\teil $g \parallel h$, g liegt oben, h unten, die Gerade k schneidet beide. An g liegt oberhalb von g rechts von k ein Winkel von $62^\circ$. Wie groß ist der Winkel an h oberhalb von h links von k?

\parallelenpaar{62}{}{}{}{}
\leerfeld °
\teil Im Dreieck ABC ist der Winkel bei C $58^\circ$ groß. Die Seiten AB und CB werden über B hinaus verlängert; zwischen den beiden Verlängerungen liegt ein Winkel von $43^\circ$. Wie groß ist der Innenwinkel β bei B? \hfill (P10 2019 OS) \\ β = \leerfeld °
\teil Im Parallelogramm ABCD ist $\alpha = 64^\circ$. Gib β, γ und δ an.

\parallelogramm[winkel={\alpha,\beta,\gamma,\delta}]{4}{2.5}{64}
β = \leerfeld °, γ = \leerfeld °, δ = \leerfeld °
\teil Die parallelen Geraden h (oben) und g (unten) werden von der Geraden k geschnitten. An h liegt unterhalb von h links von k ein Winkel von $44^\circ$; α liegt an g oberhalb von g links von k. Wie groß ist α? \hfill (P10 2015 OS) \\

\parallelenpaar{44}{}{}{}{}
α = \leerfeld °
\end{teile}
\end{aufgabe}

%% TODO Titel: Ich-kann-Satz fehlt in der Bank (Platzhalter: Kettenname) – Nr. 14
%% TODO Anweisung: ein Satz über der ersten Teilaufgabe fehlt in der Bank – Nr. 14
\begin{aufgabe}{Parallelität aus gleichen Stufenwinkeln prüfen}
\begin{teile}
\teil Die Gerade k schneidet die Geraden g (oben) und h (unten). An g liegt oberhalb von g rechts von k ein Winkel von $73^\circ$, an h oberhalb von h rechts von k ebenfalls $73^\circ$. Sind g und h parallel? \\ \kreuz{ja} \\ \kreuz{nein}
\end{teile}
\end{aufgabe}

%% TODO Titel: Ich-kann-Satz fehlt in der Bank (Platzhalter: Kettenname) – Nr. 15
%% TODO Anweisung: ein Satz über der ersten Teilaufgabe fehlt in der Bank – Nr. 15
\begin{aufgabe}{Winkel an Geraden – Fehler finden, Begründen, Anwendung}
\begin{teile}
\teil An einer Geradenkreuzung ist $\alpha = 66^\circ$. Jonas berechnet den Scheitelwinkel γ: \rechnung{\gamma &= 180^\circ - 66^\circ = 114^\circ} Finde den Fehler und gib γ richtig an.
\teil An einer Geradenkreuzung liegen α und γ einander gegenüber, β liegt dazwischen. Warum sind α und γ immer gleich groß? Begründe, ohne zu messen.
\teil Auf einem Parkplatz sind die Parkstreifen parallel. Der erste Streifen trifft den geraden Bordstein unter $62^\circ$. Welche beiden Winkel bildet der fünfte Streifen mit dem Bordstein? \leerfeld ° und \leerfeld °
\end{teile}
\end{aufgabe}
